\documentclass[a4paper,11pt]{article}
\usepackage[utf8]{inputenc}
\usepackage[T1]{fontenc}
\usepackage{geometry}
\geometry{margin=2.5cm}
\usepackage{amsmath,amssymb}

\title{QCM -- Stochastic Models and Simulation}
\author{M2 Data Science}
\date{}

\begin{document}
\maketitle

\noindent\textbf{Instructions}
\begin{itemize}
  \item Several answers may be correct for each question.
  \item Tick \textbf{all} correct answers; a partially correct selection (missing a correct option or ticking a wrong one) is marked incorrect.
  \item Duration: \textbf{20 minutes}.
\end{itemize}

\bigskip

\begin{enumerate}
\item \textbf{Poisson process and thinning.} For a homogeneous Poisson process with rate $\lambda$ on $[0,T]$:
  \begin{enumerate}
    \item The number of arrivals $N(T)$ is Poisson with mean $\lambda T$.
    \item Conditionally on $N(T)=n$, the $n$ arrival times are i.i.d.\ $\mathrm{Unif}(0,T)$.
    \item Thinning with independent Bernoulli$(p)$ marks produces a Poisson process of rate $\lambda p$.
    \item Superposition of independent Poisson processes with rates $\lambda_1,\lambda_2,\dots$ is not Poisson in general.
    \item Inter-arrival times are i.i.d.\ exponential with mean $\lambda$.
  \end{enumerate}

\item \textbf{Little’s Law and stability.} Consider a single-station queue observed in steady state.
  \begin{enumerate}
    \item Little’s Law $L=\lambda W$ holds only for $M/M/1$ queues.
    \item In any stable system, the long-run throughput always equals the \emph{external} arrival rate, even when there is blocking or abandonment.
    \item If $\lambda \ge c\mu$ in an $M/M/c$, the system is unstable (no finite steady-state mean).
    \item $L_q = \lambda W_q$ is valid even when service times are non-exponential, provided the system is stable.
    \item Little’s Law cannot be used to validate or cross-check simulation estimates because it is purely theoretical.
  \end{enumerate}

\item \textbf{$M/M/1$ intuitions and sensitivity.} For an $M/M/1$ queue with arrival rate $\lambda$ and service rate $\mu$:
  \begin{enumerate}
    \item $\rho = \lambda/\mu$ is the utilisation and must satisfy $\rho < 1$ for a stable steady state.
    \item As $\rho \to 1$, the mean waiting time $W_q$ grows approximately like $1/(1-\rho)$.
    \item Small relative errors in $\lambda$ or $\mu$ near $\rho \approx 1$ do not significantly affect $W_q$.
    \item When $\rho$ is small (e.g.\ $\rho=0.2$), $W_q$ is highly sensitive to small changes in $\lambda$.
    \item The stationary distribution of $N$ (number in system) is Poisson with mean $\rho$.
  \end{enumerate}

\item \textbf{$M/M/c$ and Erlang-C.} For an $M/M/c$ queue (Poisson arrivals, $c$ exponential servers, infinite buffer):
  \begin{enumerate}
    \item The offered load is defined as $a := \lambda/(c\mu)$, which equals the utilisation $\rho$.
    \item The Erlang-C formula gives the probability that an arrival has to wait.
    \item If $\rho$ is fixed and $c \to \infty$, the waiting probability always goes to $1$.
    \item For fixed $\lambda,\mu$, adding servers (increasing $c$) always decreases $W_q$ with roughly linear (non-diminishing) improvement.
    \item The birth--death structure uses $\lambda_n=\lambda$ and $\mu_n = \min(n,c)\mu$.
  \end{enumerate}

\item \textbf{CLT-based confidence intervals for simulation output.} You simulate $R$ independent replications of a performance metric $X$ (e.g.\ mean waiting time) and obtain sample mean $\bar X$ and sample standard deviation $s$.
  \begin{enumerate}
    \item A $95\%$ CI can be approximated by $\bar X \pm 1.96\,s/\sqrt{R}$ when $R$ is not too small.
    \item Increasing $R$ shrinks the CI width roughly like $1/\sqrt{R}$.
    \item Autocorrelation within each replication makes it impossible to apply the CLT and construct meaningful CIs.
    \item If you use common random numbers across configurations, you cannot construct valid CIs.
    \item When the distribution of $X$ is heavy-tailed, CLT-based CIs converge just as fast as in the Gaussian case and require no extra caution.
  \end{enumerate}

\item \textbf{Variance reduction and experiment design.} Consider Monte Carlo or queueing simulation experiments.
  \begin{enumerate}
    \item Antithetic variates use positively correlated samples (e.g.\ $U$ and $U'$) to reduce variance.
    \item Control variates exploit a correlated quantity with known expectation to reduce estimator variance.
    \item Using the same random number seeds for two competing system designs can reduce the variance of their \emph{difference}.
    \item Discarding a warm-up period has no real impact on bias in steady-state estimates.
    \item Stratified sampling is useless in simulation and should be avoided.
  \end{enumerate}

\item \textbf{Parameter estimation from event data.} You observe events from a Poisson process and an $M/M/1$ queue.
  \begin{enumerate}
    \item For a homogeneous Poisson process observed on $[0,T]$, the MLE for $\lambda$ is $\hat\lambda = N(T)/T$.
    \item In an $M/M/1$ with observed busy time $B(T)$ and completions $C(T)$, $\hat\mu = B(T)/C(T)$ is a natural estimator of $\mu$.
    \item If you only observe the aggregate count $N(T)$, you cannot consistently estimate $\lambda$.
    \item For $M/G/1$, estimating $\lambda$ and the first two moments of service times is enough to plug into the Pollaczek--Khinchine formula.
    \item Fisher information is only relevant for Bayesian estimators and does not help quantify precision for MLEs.
  \end{enumerate}

\item \textbf{Diagnosing and fitting from an $M/G/c$ log.} You are given an event log with \texttt{arrival\_time}, \texttt{start\_service\_time}, \texttt{completion\_time}, \texttt{service\_time}, \texttt{wait\_time}, \texttt{system\_time}, and \texttt{queue\_len\_at\_arrival} for a multi-server system.
  \begin{enumerate}
    \item By examining overlaps between \texttt{start\_service\_time} and \texttt{completion\_time}, you can infer the maximum observed number of simultaneous services and thus a lower bound on $c$.
    \item A histogram of inter-arrival times plus overdispersion of counts can reveal deviations from a homogeneous Poisson arrival process.
    \item A log-survival plot of \texttt{service\_time} can indicate a shifted exponential structure (flat region then exponential tail).
    \item If you see many arrivals with \texttt{queue\_len\_at\_arrival = 0}, the system must have only one server.
    \item If the fitted model consistently underestimates empirical waiting times, it means the model is overfitting noise and is therefore preferable.
  \end{enumerate}

\item \textbf{Model selection and structural assumptions.} You have a system where customers arrive roughly at rate $\lambda$ and are processed by several servers; some customers abandon the queue when waiting too long. You fit a simple $G/G/c$ model \emph{without} abandonment based only on completed jobs.
  \begin{enumerate}
    \item Using completions only to estimate $\lambda$ will generally underestimate the true arrival rate.
    \item Ignoring abandonment tends to make the system appear \emph{more} congested than it really is for arrivals.
    \item A no-abandonment model fitted on completions may still match the observed waiting times of served customers while completely missing the cost of lost customers.
    \item If the fitted no-abandonment model reproduces the empirical mean wait of served customers, then it is safe to use it for staffing decisions.
    \item A key diagnostic is to look for timeouts/reneging patterns (e.g.\ maximum observed waits, censored records) and decide whether a no-abandonment model is structurally acceptable.
  \end{enumerate}

\item \textbf{Designing robust simulation studies.} You need to compare two staffing policies for a queueing system using simulation, and report a recommendation with uncertainty.
  \begin{enumerate}
    \item Using different warm-up lengths and run lengths for each policy makes comparisons fairer.
    \item Common random numbers across policies reduce variance of the \emph{difference} in performance, making it easier to detect meaningful improvements.
    \item When utilisation is very high ($\rho$ close to $1$), you should run shorter simulations because queues explode and CLT approximations become unreliable.
    \item A good practice is to report both point estimates and CIs for key metrics, and to highlight when CIs for competing policies overlap strongly.
    \item If a small change in the fitted service-time distribution (e.g.\ slightly heavier tail) flips the ranking of policies, this indicates the decision is \emph{not} robust and should be reported as such.
  \end{enumerate}

\end{enumerate}

\newpage
\section*{Correction (do not distribute during exam)}

\begin{itemize}
  \item Q1: a, b, c
  \item Q2: c, d
  \item Q3: a, b
  \item Q4: b, e
  \item Q5: a, b
  \item Q6: b, c
  \item Q7: a, d
  \item Q8: a, b, c
  \item Q9: a, c, e
  \item Q10: b, d, e
\end{itemize}

\bigskip
\noindent\textbf{Scoring rule (out of 5 points)}
\begin{itemize}
  \item For each of the 50 statements (10 questions $\times$ 5 options):
    \begin{itemize}
      \item If the statement is \emph{true} and the student ticks it: $+1$ point.
      \item If the statement is \emph{true} and the student does \emph{not} tick it: $-1$ point.
      \item If the statement is \emph{false} and the student does \emph{not} tick it: $+1$ point.
      \item If the statement is \emph{false} and the student ticks it: $-1$ point.
    \end{itemize}
  \item Let $S$ be the total over the 50 decisions ($S\in[-50,50]$). The QCM mark on 5 points is
  \[
    \text{QCM\_score} \;=\; \max\bigl(0,\; \min\bigl(5,\; S/10\bigr)\bigr).
  \]
  Thus a random paper has expected score close to $0$, and only clearly better-than-random performance yields positive points.
\end{itemize}

\end{document}
